\begin{tablalibro}{3}{L{0.21}L{0.31}L{0.48}}
\caption{Tipos de datos básicos en Python}\label{tab:tipos-de-datos-basicos-en-python}\\
\toprule
\rowcolor{Papel}
\CabeceraTabla{Tipo} & \CabeceraTabla{Ejemplo} & \CabeceraTabla{Uso} \\
\midrule
\endfirsthead
\SigueTabla{3}
\toprule
\rowcolor{Papel}
\CabeceraTabla{Tipo} & \CabeceraTabla{Ejemplo} & \CabeceraTabla{Uso} \\
\midrule
\endhead
\midrule
\ContinuaTabla{3}
\endfoot
\bottomrule
\multicolumn{3}{@{}p{\linewidth}@{}}{\NotaLibro{Elaboración propia.}}\\
\endlastfoot

\Codigo{int} & \Codigo{12} & Cantidades enteras \\
\Codigo{float} & \Codigo{3.75} & Medidas aproximadas \\
\Codigo{str} & \Codigo{"Cuaderno"} & Texto y códigos \\
\Codigo{bool} & \Codigo{True} & Condiciones lógicas \\
\Codigo{NoneType} & \Codigo{None} & Ausencia de valor \\
\end{tablalibro}
